\scanpage{27}
\mspage{12}
% archive sheet number: 26
L'extension $\unsure{\mathscr E|\mathcal E}=\mathscr E_{0,3}$ de
\[
\Gamma=\Gamma_{0,3}\simeq\underbrace{\unsure{\mathfrak S_{\{0,1,\infty\}}|\mathscr S_{\{0,1,\infty\}}|G_{\{0,1,\infty\}}}}_{=\mathfrak S_3}\times\mathbb Z/2\mathbb Z
\]
\margin{engendré par $\tau$: $z\mapsto\overline z$}
par
\[
\pi=\pi_{0,3}=\{l_0,l_1,l_\infty\mid l_\infty l_1l_0=1\}=\pi_1(\underbrace{\mathbb P^1_{\mathbb C}\setminus\{0,1,\infty\}}_{\unsure{U=X_{0,3}|U=U_{0,3}}},-\overline j\}.
\]
Les classes de conjugaison \deleted{\unsure{au dessus des|au dessus de}} de scindages partiels maximaux correspondent aux points fixes de $\dot\rho,\dot\sigma_0,\dot\sigma_1,\dot\sigma_\infty$ opérant sur $X_{0,3}$, et aux comp. connexes des ens. d'invariants \inserted{(éventuelles)} de $\deleted{\dot\rho}\dot\tau,\dot\sigma_0\tau,\dot\sigma_1\tau,\dot\sigma_\infty\tau$. Les pts fixes de $\rho$ sont $P$ et $\overline P$ (dont le stabilisateur est d'ordre 6, \inserted{savoir $\varphi(\mathfrak S_3)$}) formé de $1,\dot\rho,\dot\rho^2,\dot\sigma_1\tau,\dot\sigma_2\tau,\dot\sigma_3\tau$. On a étudié le relèvement correspondant \inserted{à $P$} (de $\varphi(\mathfrak S_3)$), \unsure{on trouve} des relèvements ($\pi$-conjugués \inserted{entre eux}) relatifs à $\overline P$ s'obtiennent en \deleted{\unsure{conj}} prenant un élément de $\mathscr E$ \deleted{\unsure{et}} qui soit au dessus d'un élément de $\Gamma$ envoyant $P$ en $\overline P$ (savoir $\dot\sigma_0,\dot\sigma_1,\dot\sigma_\infty,\tau,\tau\rho,\tau\rho^2$) et en conjuguant par cet élément on peut donc prendre \deleted{$\unsure{\sigma_0}$} un des éléments $\sigma_0,\sigma_1,\sigma_\infty$ (ou $\unsure{\varepsilon_0|\epsilon_0},\varepsilon_1,\varepsilon_\infty$) ou $\tau_0,\tau_1,\tau_\infty$ ou les produits des précédents par \deleted{\unsure{un}} $\rho$ ou $\rho^2$\unsure{...}.

Le \inserted{seul} pt fixe de $\sigma_0$ est 2, dont le stabilisateur est le groupe $\simeq\mathbb Z/2\times\mathbb Z/2$ engendré par $\dot\sigma_0$ et $\tau$. On a un relèvement de ce sous-groupe
